\section{Which long-distance behavior is worth predicting?}
\label{sec:physics}

The scientific target is the evolution of a complex wave envelope, including
its phase, modulation and transfer between Fourier modes. A useful learned
model should preserve these features for unseen initial packets over a
declared propagation interval. Agreement with the same Dysthe equation tests
the surrogate; agreement with measured water waves tests the physical model.
These are separate comparisons.

\subsection{Modulation, focusing and recurrence}

Dysthe's 1979 calculation identifies induced mean flow as a central
correction to modulation-instability predictions \cite{dysthe1979}.
It motivates measuring early sideband growth and the onset of nonlinear
exchange. Its temporal potential-envelope formulas must not be substituted
directly into the present spatial free-surface equation. Instead, use the
linearization of \eqref{eq:baseline} and compare small perturbations of a
carrier at several modulation frequencies, amplitudes and relative phases.
A growing perturbation is evidence of instability of that carrier, not
evidence of wave breaking or loss of smoothness.

Localized packets offer a complementary test. Lo and Mei studied asymmetric
leaning and splitting, together with unequal sideband growth, using a
higher-order envelope model and comparisons with experiments \cite{lomei}.
For Gaussian and sech-like initial packets, measure maximum amplification,
the distance and profile location of that maximum, packet asymmetry and
spectral broadening. Varying width, chirp and amplitude tests distinct
interactions between dispersion and nonlinearity. A small average field
error can still hide an incorrect focusing distance or an underestimated
peak, so these observables require separate error measures.

Recurrence tests a longer sequence of transfers from carrier to sidebands
and back. Armaroli, Brunetti and Kasparian compare full surface-envelope
HONLS evolution with a three-mode reduction and distinguish recurrence
classes through relative phase \cite{recurrence}. Their reduction is useful
as a cheap comparator only while omitted modes remain small; their
normalization requires an explicit conversion before quantitative reuse.
Measure carrier depletion, sideband imbalance, recurrence distance and
the complex field at return. A recurrence defect may minimize over a constant
phase and a profile translation, but both fitted shifts and the unaligned
error should also be reported. Repeated maxima alone do not establish a
breather or accurate recurrence.

Two-packet interactions provide an optional harder test. Fedele and Dutykh
numerically observe radiation and changed profiles after spatial-Dysthe
collisions, contrasting with their NLS control \cite{fedele}. Compare the
joint evolution with separately propagated packets and quantify outgoing
shape, phase, displacement and radiation. The comparison must resolve the
small interaction signal and account for overlap; a subtraction artifact
is not emitted radiation, and a collision calculation is not a proof of
nonintegrability.

\subsection{Spectral centroid and spectral peak are different observables}

For the Fourier convention
$u(\xi,\tau)=\sum_n\widehat u_n(\xi)e^{ik_n\tau}$,
$k_n=2\pi n/L$, define
\[
 \overline k=\frac{\sum_n k_n\abs{\widehat u_n}^2}
                     {\sum_n\abs{\widehat u_n}^2},
 \qquad
 k_{\mathrm{peak}}\in\operatorname*{arg\,max}_{k_n}
                         \abs{\widehat u_n}^2.
\]
The centroid averages all modes; the peak can jump when two modal powers
cross. Armaroli et al. show why a temporary mean upshift can coexist with
a peak downshift \cite{recurrence}. Neither statistic should substitute for
the other.

The following identities are derived directly for our selected baseline.
All integrals below are over the periodic profile circle $\TT$ of length $L$.

\begin{proposition}[Centroid constraint for the continuum equation]
\label{prop:centroid}
Let $u$ be a sufficiently smooth periodic solution of \eqref{eq:baseline}
with real $\epsilon$, and define
\[
 A=\int_{\TT}|u|^2\,\dd\tau,\qquad
 P=\operatorname{Im}\int_{\TT}\overline u\,u_\tau\,\dd\tau,\qquad
 Q=\int_{\TT}|u|^4\,\dd\tau.
\]
Then $A$ and $P-\epsilon Q$ are constant. If $A>0$, then
\begin{equation}
 \overline k(\xi)-\overline k(0)
 =\frac{\epsilon}{A}\bigl(Q(\xi)-Q(0)\bigr).
 \label{eq:centroid-constraint}
\end{equation}
\end{proposition}
\begin{proof}
Set $\rho=|u|^2$ and
$J=\operatorname{Im}\int_{\TT}\overline u^{\,2}u_\tau^2\,\dd\tau$.
Differentiation and periodic integration by parts give
\[
 A'=-10\epsilon\int_{\TT}\rho\rho_\tau\,\dd\tau=0,
 \qquad P'=-4\epsilon J,\qquad Q'=-4J.
\]
The mean-flow contribution to $P'$ vanishes because
$\int_{\TT}\rho_\tau\D\rho\,\dd\tau=0$: $\D$ is self-adjoint,
commutes with $\partial_\tau$, and preserves real functions.
Thus $(P-\epsilon Q)'=0$. Parseval's identity gives
$\overline k=P/A$, proving \eqref{eq:centroid-constraint}.
\end{proof}

This relation links centroid motion to changes in concentration; it does
not prescribe a universal direction of drift. It is a continuum diagnostic,
so compare a computed violation with the refined solver's own defect.
Projection and time discretization can affect invariant preservation.

Here $k$ is an \emph{envelope} frequency conjugate to $\tau$. With the stated
carrier convention $e^{i(x-t)}$ and $\tau=\epsilon(2x-t)$, a mode has
physical angular frequency $\omega_0(1+\epsilon k)$ in the positive-carrier
band. Changing the Fourier or carrier convention changes this translation.
The full reconstructed surface includes bound harmonics; its spectral mean
need not equal the first-harmonic envelope statistic.

\subsection{What a long calculation can and cannot show}

Long propagation means large $\xi$. Packet width and modulation bandwidth
are distinct diagnostics; chirp can broaden the spectrum. The carrier
remains a deep-water wave in this model. Moreover,
$\xi=\epsilon^2 k_0X$ means that the same dimensionless distance corresponds
to different physical distances $X$ when steepness varies.

Periodicity also changes the interpretation of long runs. Since $A$ is
conserved,
\begin{equation}
 \norm{u(\xi,\cdot)}_{L^\infty(\TT)}
 \geq\sqrt{A/L}.
 \label{eq:periodic-floor}
\end{equation}
Indeed, $A\leq L\norm{u}_{L^\infty}^2$. A fixed periodic box therefore
cannot demonstrate decay of a nonzero packet to zero. Radiation can return
through the boundary; localized-packet conclusions require domain checks
and examination of the exterior field.

Track maximum effective steepness and occupied bandwidth as well as
numerical resolution. For a conservative pilot screen, one could require
$\epsilon\max|u|\leq.1$ and $\epsilon k_{95}\leq.1$, where $k_{95}$
is the smallest nonnegative cutoff containing $95\%$ of the envelope
action in $|k|\leq k_{95}$. Record the full tails as well. These proposed
screening numbers do not certify asymptotic validity and must be
reconsidered for a measured wave population. Trulsen and Dysthe's
broader-band extension illustrates
that relaxing bandwidth restrictions requires a changed model
\cite{broadband}. No universal amplitude cutoff or indefinitely valid
propagation distance follows from a successful envelope computation.

Physical long-distance comparisons need further care. Carter and Govan
derive a viscous Dysthe model and compare downshift observables with tank
data \cite{viscous}. A later comparison across four experiments finds that
model performance depends on the observable and initial sideband seeding;
no model captures every measured evolution \cite{experiments}.
Transverse modulation can also transfer energy and affect downshift in a
wide tank \cite{transverse}. The present inviscid, one-profile equation
does not represent dissipation, transverse evolution or breaking.

Finally, accurate finite-distance fields and bounded sampled derivatives
support a resolved computational observation. They do not establish global
smooth existence. A small training residual does not establish that either:
causal PINN training addresses optimization and information propagation
\cite{causal}, while existence and uniform error certification require
additional mathematical estimates.
